\documentclass[11pt]{article}
\usepackage{amsmath, amssymb}
\usepackage[margin=1in]{geometry}
\setlength{\parskip}{0.5em}

\title{Identified Combinations in the Latent Exposure Model}
\author{Johnson Vo}
\date{}

\begin{document}
\maketitle

Working notes. The setting follows Gustafson (2010), ``Bayesian Inference for Partially Identified Models'': a transparent reparameterization $\phi = (\phi_I, \phi_N)$ where the data depend only on $\phi_I$. These are derived by hand and the second-moment argument in Section 4 has not been checked numerically.

\section{Model}

\begin{align*}
A_i &= \gamma_0 + X_i^\top \gamma_A + e_i, & e_i &\sim N(0, \sigma_A^2),\\
Z_i &= \beta_{Z1} A_i + \beta_{Z2} D_i + \varepsilon_i, & \varepsilon_i &\sim N(0, \sigma_Z^2),\\
Y_i &\sim \text{Poisson}\big(N_i \exp(\beta_0 + \beta_A A_i + X_i^\top \beta_X + \phi_i)\big), & \phi &\sim \text{CAR}(\sigma_\phi, \rho).
\end{align*}

Only $(Y, Z, X, D, N)$ are observed. The anchored model sets $\beta_{Z1} = 1$.

\section{Surrogate alone}

Marginally,
\[
Z_i \mid X_i, D_i \sim N\big(\beta_{Z1}\gamma_0 + \beta_{Z1} X_i^\top \gamma_A + \beta_{Z2} D_i,\; \beta_{Z1}^2 \sigma_A^2 + \sigma_Z^2\big).
\]
So the data identify $\alpha_0 = \beta_{Z1}\gamma_0$, $\alpha_X = \beta_{Z1}\gamma_A$, $\beta_{Z2}$, and the total variance $v = \beta_{Z1}^2\sigma_A^2 + \sigma_Z^2$.

\section{Scale ridge}

Let $U_i = \beta_{Z1} A_i$ (exposure on the scale of $Z$). Then
\[
Z_i = U_i + \beta_{Z2} D_i + \varepsilon_i, \qquad U_i \sim N(\alpha_0 + X_i^\top \alpha_X,\; s_U^2), \qquad \beta_A A_i = \kappa U_i,
\]
with
\[
s_U = \beta_{Z1}\sigma_A, \qquad \kappa = \beta_A / \beta_{Z1}.
\]
The likelihood depends on the exposure and error parameters only through $(\alpha_0, \alpha_X, \beta_{Z2}, s_U, \sigma_Z, \kappa)$, plus $\beta_0$, $\beta_X$, $\sigma_\phi$, $\rho$. It does not depend on the rescaling
\[
\beta_{Z1} \to \beta_{Z1}/c, \qquad \sigma_A \to c\,\sigma_A, \qquad \beta_A \to \beta_A/c,
\]
which leaves $s_U$ and $\kappa$ unchanged. So $\phi_N = \beta_{Z1}$ and everything else is in $\phi_I$.

Check against the free-slope sparse run (100 areas, 4 chains). Posterior means by chain:

\begin{center}
\begin{tabular}{lcccc}
 & chain 1 & chain 2 & chain 3 & chain 4 \\
$\beta_{Z1}$ & 1.03 & 1.02 & 1.71 & 1.53 \\
$\sigma_A$ & 1.67 & 1.80 & 1.00 & 1.16 \\
$\sigma_A \beta_{Z1}$ & 1.61 & 1.64 & 1.64 & 1.67 \\
$\beta_A / \beta_{Z1}$ & 0.115 & 0.110 & 0.117 & 0.111 \\
\end{tabular}
\end{center}

Truth: $\sigma_A\beta_{Z1} = 1.5$, $\beta_A/\beta_{Z1} = 0.1$.

\section{Reliability (with the scale fixed)}

Let $\lambda = s_U^2/(s_U^2 + \sigma_Z^2)$ be the reliability, so $\sigma_Z^2 = (1-\lambda)v$ and $s_U^2 = \lambda v$. Using the normal posterior of $U_i$ given $Z_i$,
\[
\log E[Y_i \mid Z_i, X_i, D_i] = \beta_0 + X_i^\top\beta_X + \kappa(1-\lambda)(\alpha_0 + X_i^\top\alpha_X) + \kappa\lambda (Z_i - \beta_{Z2} D_i) + \tfrac{1}{2}\kappa^2 v_{\text{post}} + \phi_i,
\]
where $v_{\text{post}} = s_U^2\sigma_Z^2/(s_U^2+\sigma_Z^2)$.

The slope on $Z$ identifies only $c_Z = \kappa\lambda$. The coefficient on $D$ is $-c_Z\beta_{Z2}$, which adds nothing new. The mean function therefore cannot separate $\kappa$ from $\lambda$: the usual need for the error variance from outside information.

The only other information on $\sigma_Z$ is the extra variance $\kappa^2 v_{\text{post}}$ that the noise in $U$ adds to the counts. This is confounded with the variance of $\phi$. They can be separated only through spatial structure ($\phi$ is spatially correlated, the noise in $U$ is not). This is a hypothesis; it would explain why $\sigma_Z$ is weakly identified and why $\rho$ has a wide posterior. Not verified.

\section{Consequences for the estimand}

\begin{itemize}
\item The shape of the exposure--response curve is identified in terms of $U = \beta_{Z1}A$, but not its horizontal scale. Our grid $a$ is in units of $A$, so $\theta(a)$ depends on $\beta_{Z1}$. Anchoring $\beta_{Z1} = 1$ is the assumption that $Z$ is on the same scale as $A$.
\item With the scale fixed, $\kappa = c_Z/\lambda$ and $\lambda = 1 - \sigma_Z^2/v \in (0, 1]$. This gives $\kappa \ge c_Z$ (attenuation) and no upper bound unless $\sigma_A$ is bounded below. With $\sigma_A \ge 0.5$ and $v \approx 3.25$: $\lambda \ge 0.25/3.25 \approx 0.077$, so $\kappa \le 13\,c_Z$. This ignores the second-moment information above, so the true region may be narrower.
\item $\sigma_Z$ is confined to $(0, \sqrt{v})$, about $(0, 1.8)$ here. A prior on $\sigma_Z$ with mass outside this region (for example prior mean 2) is pulled to the upper edge, which fits the posterior mean of 1.24 seen in the prior sensitivity run.
\item By Gustafson (2010), the posterior of $\phi_N$ given $\phi_I$ equals its prior conditional however large $n$ is, so the posterior mean of the target converges to a prior-dependent quantity.
\end{itemize}

Rough numbers for the simulation ($s_U = 1.5$, $\sigma_Z = 1$, $\kappa = 0.1$): $v = 3.25$, $\lambda \approx 0.69$, $c_Z \approx 0.069$.

\end{document}
